\section{DPO: Direct Preference Optimization}

\subsection{Cómo prescindir del modelo de recompensa}

\begin{importantbox}{La idea central de DPO}
El problema de RL con restricción KL que plantea RLHF tiene una \textbf{solución de forma cerrada}:
$$\pi^*(y \mid x) = \frac{1}{Z(x)} \pi_{\mathrm{ref}}(y \mid x) \exp\left(\frac{1}{\beta} r(x, y)\right)$$

A la inversa, la recompensa puede expresarse como función de la política:
$$r(x, y) = \beta \log \frac{\pi^*(y \mid x)}{\pi_{\mathrm{ref}}(y \mid x)} + \beta \log Z(x)$$

Al sustituir esta expresión en el modelo de Bradley-Terry, el término $Z(x)$ se cancela y se obtiene la pérdida de DPO:
$$\mathcal{L}_{\mathrm{DPO}}(\theta) = -\mathbb{E}_{(x, y_w, y_l)}\left[\log \sigma\left(\beta \log \frac{\pi_\theta(y_w \mid x)}{\pi_{\mathrm{ref}}(y_w \mid x)} - \beta \log \frac{\pi_\theta(y_l \mid x)}{\pi_{\mathrm{ref}}(y_l \mid x)}\right)\right]$$
\end{importantbox}

Ventajas de DPO:
\begin{itemize}
    \item No requiere entrenar un modelo de recompensa por separado.
    \item No requiere entrenamiento con RL (no hace falta PPO).
    \item Basta con un ciclo estándar de entrenamiento supervisado.
    \item El entrenamiento es más estable y más sencillo.
\end{itemize}

\subsection{La intuición detrás de DPO}

La pérdida de DPO hace implícitamente dos cosas:
\begin{enumerate}
    \item \textbf{Aumenta} la probabilidad de la respuesta preferida $y_w$.
    \item \textbf{Reduce} la probabilidad de la respuesta no preferida $y_l$.
\end{enumerate}

El peso se ajusta automáticamente mediante la «diferencia de recompensa implícita» dentro de la función $\sigma$.

\begin{knowledgebox}{DPO vs. RLHF}
\begin{itemize}
    \item \textbf{RLHF}: se entrena un modelo de recompensa $\to$ optimización con RL. Son dos etapas y se necesita infraestructura de RL.
    \item \textbf{DPO}: optimiza la política directamente sobre los datos de preferencias. Es una sola etapa y basta con el código de SFT.
    \item En teoría, DPO y RLHF optimizan \textbf{el mismo objetivo}.
    \item En la práctica, DPO es más sencillo, pero en algunos escenarios puede ser menos flexible que RLHF.
\end{itemize}
\end{knowledgebox}

\subsection{Práctica y elección de parámetros}

Elementos clave durante el entrenamiento con DPO:
\begin{itemize}
    \item Elegir una reference policy estable; por lo general se usa un instruction finetuned checkpoint.
    \item Ajustar $\beta$ para controlar el grado de preferencia: cuanto menor es $\beta$, más se acerca la política a la reference; cuanto mayor es, más fácil resulta reforzar la preferencia.
    \item Los pesos de muestreo deben equilibrar la cabeza y la cola de la distribución, para evitar que el reward model prefiera únicamente las respuestas de alta probabilidad.
\end{itemize}

\begin{table}[H]
\centering
\begin{tabular}{p{0.32\textwidth} p{0.58\textwidth}}
\toprule
Hiperparámetro & Configuración recomendada \\
\midrule
$\beta$ & 0.1-1.0; hacer un sweep según la intensidad de la preferencia y observar el trade-off entre KL y accuracy \\
Reference policy & SFT checkpoint + small RL finetuning for stability \\
Composición del lote & Mantener una balanced ratio de pares winning/losing para evitar el collapse \\
\bottomrule
\end{tabular}
\caption{Ajuste de hiperparámetros y práctica de entrenamiento en DPO}
\end{table}

\begin{knowledgebox}{DPO Training Insight}
DPO convierte el reward scoring en una diferencia de logits; por eso, al observar el shift de $\log \pi_\theta / \pi_{\mathrm{ref}}$, puede juzgarse de forma intuitiva cuánto se ha reforzado la preferencia.
\end{knowledgebox}

\subsection{Resumen de la sección}

DPO sustituye la solución de forma cerrada del objetivo de RL en el modelo de Bradley-Terry y así reduce la optimización de preferencias a un problema de aprendizaje puramente supervisado. Combinar un sweep de $\beta$ con la elección de la política de referencia permite equilibrar la precisión y la estabilidad.
